La Tabla~\ref{tab:parametros} separa tres escenarios que después van a ser tres
curvas distintas en el mismo gráfico. El único cambio entre el primero y el
segundo es $R_L$, la resistencia del bobinado: no mueve la frecuencia natural,
porque $\omega_0$ no depende de la resistencia, pero se lleva puesto un
\SI{10}{\percent} del factor de mérito.

\begin{table}[H]
\centering
\caption{Parámetros del filtro en los tres escenarios que compara el informe.}
\label{tab:parametros}
\small
\begin{tabular}{lrrrr}
\toprule
\textbf{Escenario} & $R_t$ & $\fnat$ & $Q$ & $|H|_{\max}$ \\
\midrule
Teórico ideal ($R_L = 0$)      & \SI{330}{\ohm}   & \SI{1073.0}{\hertz} & 2,043 & \SI{6.47}{\decibel} \\
Simulado (con $R_L$ nominal)   & \SI{368}{\ohm}   & \SI{1073.0}{\hertz} & 1,832 & \SI{5.59}{\decibel} \\
Medido (componentes reales)    & \SI{373.4}{\ohm} & \SI{1077.7}{\hertz} & 1,750 & \SI{5.23}{\decibel} \\
\bottomrule
\end{tabular}
\end{table}

\noindent
Con $\alpha = R_t/2L = \SI{1650}{\per\second}$ y
$\omega_d = \SI{6537}{\radian\per\second}$ en el caso ideal, la
Ecuación~\eqref{eq:escalon} anticipa un sobrepico de
$M_p = e^{-\pi\alpha/\omega_d} = \SI{45}{\percent}$ y una oscilación amortiguada
de \SI{1040}{\hertz}, que es lo que se va a ver en el osciloscopio.
